\section*{第 2 章}
\addcontentsline{toc}{section}{第 2 章}

\exercise{习题 2.1}
\textbf{题目：}
设 $m(t)$ 是功率为 $2\,\mathrm{W}$ 的实信号，求
\[
z(t)=m(t)+j\cdot 2m(t)
\]
的功率。

\par\textbf{解：}
复信号的功率等于实部功率与虚部功率之和，因此
\[
P_z=2+4\times 2=10\,\mathrm{W}.
\]
也可写成
\[
z(t)=m(t)(1+2j),
\qquad
\lvert z(t)\rvert^2=m^2(t)(1+2^2)=5m^2(t).
\]
按时间取平均，因 $\overline{m^2(t)}=2\,\mathrm{W}$，故
\[
\overline{\lvert z(t)\rvert^2}=5\times 2=10\,\mathrm{W}.
\]

\medskip
\exercise{习题 2.2}
\textbf{题目：}
考虑 4 个信号
\[
w(t)=
\begin{cases}
1,&0\le t<2T,\\
0,&\text{其他},
\end{cases}
\quad
x(t)=
\begin{cases}
1,&0\le t<T,\\
0,&\text{其他},
\end{cases}
\]
\[
y(t)=
\begin{cases}
1,&T\le t<2T,\\
0,&\text{其他},
\end{cases}
\quad
z(t)=
\begin{cases}
1,&0\le t<T,\\
-\dfrac12,&T\le t<2T,\\
0,&\text{其他},
\end{cases}
\]
求它们两两之间的内积。

\par\textbf{解：}
\[
\int_{-\infty}^{\infty}x(t)w(t)\,dt
=\int_{-\infty}^{\infty}x(t)\,dt=T,
\]
\[
\int_{-\infty}^{\infty}w(t)y(t)\,dt
=\int_{-\infty}^{\infty}y(t)\,dt=T,
\]
\[
\int_{-\infty}^{\infty}w(t)z(t)\,dt
=\int_{-\infty}^{\infty}z(t)\,dt=\frac{T}{2},
\]
\[
\int_{-\infty}^{\infty}x(t)y(t)\,dt=0,
\qquad
\int_{-\infty}^{\infty}x(t)z(t)\,dt
=\int_{-\infty}^{\infty}x(t)\,dt=T,
\]
\[
\int_{-\infty}^{\infty}y(t)z(t)\,dt
=-\frac12\int_{-\infty}^{\infty}y(t)\,dt
=-\frac{T}{2}.
\]
全部结果可写成矩阵
\[
\begin{pmatrix}
2T&T&T&\dfrac{T}{2}\\[4pt]
T&T&0&T\\[4pt]
T&0&T&-\dfrac{T}{2}\\[4pt]
\dfrac{T}{2}&T&-\dfrac{T}{2}&\dfrac{5T}{4}
\end{pmatrix},
\]
其中第 $1,2,3,4$ 行及第 $1,2,3,4$ 列分别对应信号 $w,x,y,z$。

\medskip
\exercise{习题 2.3}
\textbf{题目：}
证明
\[
\delta(t-t_1)\delta(t-t_2)=\delta(t-t_1)\delta(t_1-t_2).
\]

\par\textbf{证明：}
将冲激函数看成极限
\[
\delta(x)=\lim_{X\to 0}\left\{\frac{1}{X}\sinc\!\left(\frac{x}{X}\right)\right\}.
\]
则
\[
\delta(t-t_1)\delta(t-t_2)
=\delta(t-t_1)\lim_{T_2\to 0}\left\{\frac{1}{T_2}\sinc\!\left(\frac{t-t_2}{T_2}\right)\right\}.
\]
由冲激函数的抽样性质，
\[
\delta(t-t_1)\cdot \frac{1}{T_2}\sinc\!\left(\frac{t-t_2}{T_2}\right)
=\delta(t-t_1)\cdot \frac{1}{T_2}\sinc\!\left(\frac{t_1-t_2}{T_2}\right).
\]
因此
\[
\delta(t-t_1)\delta(t-t_2)
=\delta(t-t_1)\lim_{T_2\to 0}\left\{\frac{1}{T_2}\sinc\!\left(\frac{t_1-t_2}{T_2}\right)\right\}
=\delta(t-t_1)\delta(t_1-t_2).
\]

\medskip
\exercise{习题 2.4}
\textbf{题目：}
求信号
\[
s(t)=Ae^{j(2\pi f_0 t+\varphi)}
\]
的傅氏变换与功率谱密度。

\par\textbf{解：}
写成
\[
s(t)=Ae^{j\varphi}e^{j2\pi f_0 t}.
\]
这是复幅度为 $Ae^{j\varphi}$、频率为 $f_0$ 的复单频信号，因此
\[
S(f)=Ae^{j\varphi}\delta(f-f_0).
\]
功率谱密度为
\[
P_s(f)=A^2\delta(f-f_0).
\]

\medskip
\exercise{习题 2.5}
\textbf{题目：}
求信号
\[
g(t)=\rect\!\left(\frac{t}{T_s}\right)\cos\!\left(\frac{\pi t}{T_s}\right)
\]
的能量、傅氏变换、能量谱密度与主瓣带宽。

\par\textbf{解：}
其能量为
\[
E_g=\int_{-\infty}^{\infty}g^2(t)\,dt
=\int_{-T_s/2}^{T_s/2}\cos^2\!\left(\frac{\pi t}{T_s}\right)\,dt
=\frac{T_s}{2}.
\]
又有
\[
g(t)=\rect\!\left(\frac{t}{T_s}\right)\cdot \frac12
\left(e^{j\pi t/T_s}+e^{-j\pi t/T_s}\right).
\]
由于
\[
\mathcal{F}\left\{\rect\!\left(\frac{t}{T_s}\right)\right\}
=T_s\sinc(fT_s),
\]
乘以复指数相当于频移，因此
\[
G(f)=\frac{T_s}{2}\left\{
\sinc\!\left[\left(f-\frac{1}{2T_s}\right)T_s\right]
\;+\;
\sinc\!\left[\left(f+\frac{1}{2T_s}\right)T_s\right]
\right\}.
\]
化简得
\[
G(f)=\frac{2T_s\cos(\pi fT_s)}{\pi(1-4f^2T_s^2)}.
\]
能量谱密度为
\[
\lvert G(f)\rvert^2
=\frac{4T_s^2\cos^2(\pi fT_s)}{\pi^2(1-4f^2T_s^2)^2}.
\]
能量谱密度的第一个零点不在 $f=\dfrac{1}{2T_s}$ 处，因为该点是可去奇点；第一个真正零点满足
\[
\pi fT_s=\frac{3\pi}{2},
\]
故
\[
f=\frac{3}{2T_s}.
\]
因此主瓣带宽为
\[
\frac{3}{2T_s}.
\]

\medskip
\exercise{习题 2.6}
\textbf{题目：}
求信号
\[
\sinc\!\left(\frac{t}{T}\right)
\]
的能量。

\par\textbf{解：}
\[
\mathcal{F}\left\{\sinc\!\left(\frac{t}{T}\right)\right\}
=T\rect(fT).
\]
因此能量谱密度为
\[
\lvert S(f)\rvert^2=T^2\rect(fT).
\]
由帕塞瓦尔定理，能量为
\[
E=\int_{-\infty}^{\infty}\lvert S(f)\rvert^2\,df
=T^2\cdot \frac{1}{T}=T.
\]

\medskip
\exercise{习题 2.7}
\textbf{题目：}
证明：若一个系统满足线性、时不变这两个特性，则其输入为
\[
e^{j2\pi ft}
\]
时，输出为
\[
c\cdot e^{j2\pi ft},
\]
其中 $c$ 是复系数。

\par\textbf{证明：}
设输入 $e^{j2\pi ft}$ 时输出为 $y(t)$。由时不变性可知，对任意实数 $t_0$，输入
\[
e^{j2\pi f(t+t_0)}
\]
时，输出为 $y(t+t_0)$。另一方面，
\[
e^{j2\pi f(t+t_0)}=e^{j2\pi ft_0}e^{j2\pi ft},
\]
根据线性特性，对应输出应为
\[
e^{j2\pi ft_0}y(t).
\]
因此
\[
y(t+t_0)=e^{j2\pi ft_0}y(t).
\]
令 $t=0$，记 $c=y(0)$，则
\[
y(t_0)=ce^{j2\pi ft_0}.
\]
由于 $t_0$ 任意，故
\[
y(t)=ce^{j2\pi ft}.
\]

\medskip
\exercise{习题 2.8}
\textbf{题目：}
将周期信号
\[
s(t)=\sum_{n=-\infty}^{\infty}\delta(t-nT)
\]
通过一个传递函数为 $H(f)$、冲激响应为 $h(t)$ 的滤波器，求输出信号 $y(t)$ 的傅氏变换、傅里叶级数展开式与功率谱密度。

\par\textbf{解：}
每个冲激通过系统后的输出为 $h(t-nT)$，因此
\[
y(t)=\sum_{n=-\infty}^{\infty}h(t-nT).
\]
输入信号的傅氏变换为
\[
S(f)=\frac{1}{T}\sum_{m=-\infty}^{\infty}\delta\!\left(f-\frac{m}{T}\right),
\]
故输出信号的傅氏变换为
\[
Y(f)=H(f)S(f)
=\frac{1}{T}\sum_{m=-\infty}^{\infty}H\!\left(\frac{m}{T}\right)\delta\!\left(f-\frac{m}{T}\right).
\]
作傅氏反变换得
\[
y(t)=\frac{1}{T}\sum_{m=-\infty}^{\infty}H\!\left(\frac{m}{T}\right)e^{j2\pi mt/T},
\]
这就是输出信号的傅里叶级数展开式。

各离散谱线的功率谱密度相加得
\[
P_y(f)=\frac{1}{T^2}\sum_{m=-\infty}^{\infty}
\left|H\!\left(\frac{m}{T}\right)\right|^2
\delta\!\left(f-\frac{m}{T}\right).
\]
因此输出功率为
\[
P_y=\int_{-\infty}^{\infty}P_y(f)\,df
=\frac{1}{T^2}\sum_{m=-\infty}^{\infty}
\left|H\!\left(\frac{m}{T}\right)\right|^2.
\]

\medskip
\exercise{习题 2.9}
\textbf{题目：}
求矩形脉冲
\[
s(t)=\rect\!\left(\frac{t}{T}\right)
\]
的自相关函数。

\par\textbf{解：}
按定义
\[
R_s(\tau)=\int_{-\infty}^{\infty}s(t+\tau)s(t)\,dt
=\int_{-\infty}^{\infty}\rect\!\left(\frac{t+\tau}{T}\right)\rect\!\left(\frac{t}{T}\right)\,dt.
\]
当 $\lvert\tau\rvert\ge T$ 时，两矩形不重叠，故 $R_s(\tau)=0$；当 $\lvert\tau\rvert<T$ 时，重叠区间宽度为 $T-\lvert\tau\rvert$，故
\[
R_s(\tau)=
\begin{cases}
T-\lvert\tau\rvert,&\lvert\tau\rvert\le T,\\
0,&\text{其他}.
\end{cases}
\]
也可由
\[
S(f)=T\cdot \sinc(fT),\qquad
\lvert S(f)\rvert^2=T^2\sinc^2(fT)
\]
再作反变换得到同样结论。

\medskip
\exercise{习题 2.10}
\textbf{题目：}
能量为 $E_s$ 的信号 $s(t)$ 通过一个冲激响应为 $h(t)$ 的线性时不变系统后，在 $t=T$ 时刻采样得到采样值 $z$。若已知 $h(t)$ 的能量为 $E_h$，求能使 $z$ 最大的 $h(t)$。

\par\textbf{解：}
采样值为
\[
z=y(T)=\int_{-\infty}^{\infty}h(\tau)s(T-\tau)\,d\tau
=\int_{-\infty}^{\infty}h(t)s(T-t)\,dt.
\]
这就是 $h(t)$ 与 $s(T-t)$ 的内积。由柯西-施瓦茨不等式，内积达到最大时应有
\[
h(t)=K\,s(T-t).
\]
又因题设要求
\[
E_h=\int_{-\infty}^{\infty}\lvert h(t)\rvert^2\,dt
=K^2E_s,
\]
故
\[
K=\sqrt{\frac{E_h}{E_s}}.
\]
因此最优冲激响应为
\[
h(t)=\sqrt{\frac{E_h}{E_s}}\,s(T-t).
\]

\medskip
\exercise{习题 2.11}
\textbf{题目：}
设有带通信号
\[
s(t)=a(t)\cos(2\pi f_c t)-b(t)\sin(2\pi f_c t),
\]
其中 $a(t),b(t)$ 是基带信号。求 $s(t)$ 以 $\cos(2\pi f_c t+\theta)$ 为参考载波的复包络。

\par\textbf{解：}
所求复包络 $s_L(t)$ 满足
\[
s(t)=\Re\left\{s_L(t)e^{j(2\pi f_c t+\theta)}\right\}.
\]
又因为
\[
s(t)=\Re\left\{[a(t)+jb(t)]e^{j2\pi f_c t}\right\}
=\Re\left\{[a(t)+jb(t)]e^{-j\theta}e^{j(2\pi f_c t+\theta)}\right\},
\]
因此
\[
s_L(t)=[a(t)+jb(t)]e^{-j\theta}.
\]

\medskip
\exercise{习题 2.12}
\textbf{题目：}
设有窄带信号
\[
x(t)=m(t)\cos(2\pi f_c t+\varphi),
\]
其中 $\varphi$ 是定值，$m(t)$ 是实基带信号，且其带宽远小于 $f_c$。求 $x(t)$ 的复包络 $x_L(t)$、希尔伯特变换 $\hat{x}(t)$，以及 $\hat{x}(t)$ 的复包络。

\par\textbf{解：}
将 $x(t)$ 写成
\[
x(t)=\Re\left\{m(t)e^{j(2\pi f_c t+\varphi)}\right\}
=\Re\left\{m(t)e^{j\varphi}e^{j2\pi f_c t}\right\}.
\]
若默认参考载波为 $\cos(2\pi f_c t)$，则
\[
x_L(t)=m(t)e^{j\varphi}.
\]
由于
\[
m(t)e^{j(2\pi f_c t+\varphi)}
\]
只有正频率分量，所以它就是 $x(t)+j\hat{x}(t)$ 对应的解析信号，从而
\[
\hat{x}(t)=\Im\left\{m(t)e^{j(2\pi f_c t+\varphi)}\right\}
=m(t)\sin(2\pi f_c t+\varphi).
\]
再写成
\[
\hat{x}(t)=\Re\left\{-jm(t)e^{j\varphi}e^{j2\pi f_c t}\right\},
\]
故 $\hat{x}(t)$ 的复包络为
\[
\hat{x}_L(t)=-jm(t)e^{j\varphi}=-jx_L(t).
\]

\medskip
\exercise{习题 2.13}
\textbf{题目：}
设有窄带信号
\[
x(t)=I(t)\cos(2\pi f_c t+\varphi)-Q(t)\sin(2\pi f_c t+\varphi),
\]
其中 $I(t),Q(t)$ 是实基带信号，且其带宽远小于 $f_c$。

求 $x(t)$ 的希尔伯特变换 $\hat{x}(t)$，并分别以 $\cos(2\pi f_c t+\varphi)$ 和 $\cos(2\pi f_c t)$ 为参考载波，求 $x(t),\hat{x}(t)$ 的复包络。

\par\textbf{解：}
由关系
\[
\Re\{(a+jb)e^{jc}\}=a\cos c-b\sin c
\]
可知
\[
x(t)=\Re\left\{[I(t)+jQ(t)]e^{j(2\pi f_c t+\varphi)}\right\}.
\]
对应解析信号的虚部即为希尔伯特变换，因此
\[
\hat{x}(t)=\Im\left\{[I(t)+jQ(t)]e^{j(2\pi f_c t+\varphi)}\right\}
=Q(t)\cos(2\pi f_c t+\varphi)+I(t)\sin(2\pi f_c t+\varphi).
\]

以 $\cos(2\pi f_c t+\varphi)$ 为参考载波时，
\[
x_{L,\varphi}(t)=I(t)+jQ(t),
\qquad
\hat{x}_{L,\varphi}(t)=Q(t)-jI(t)=-j\,x_{L,\varphi}(t).
\]

以 $\cos(2\pi f_c t)$ 为参考载波时，
\[
x_L(t)=[I(t)+jQ(t)]e^{j\varphi},
\qquad
\hat{x}_L(t)=[Q(t)-jI(t)]e^{j\varphi}=-j\,x_L(t).
\]
